\documentclass{svproc}

\usepackage[utf8]{inputenc}
\usepackage{url}

\usepackage{graphicx}
\usepackage{amsmath,amssymb,amsfonts}
\usepackage{algorithm}
\usepackage{algpseudocode}
\usepackage{booktabs}
\usepackage{hyperref}
\usepackage{cite}

\begin{document}

\title{Stationary-Coupled Modular Boolean Networks:\\Formalization, Complexity Bounds, and High-Performance C++ Engine}

\titlerunning{Stationary-Coupled Modular Boolean Networks}

\author{Carlos R. P. Tovar\inst{1} \and David C. Martins-Jr\inst{1} \and Luiz C. S. Rozante\inst{1}}

\authorrunning{C. R. P. Tovar et al.}

\institute{Center for Mathematics, Computing and Cognition, Federal University of ABC (UFABC), Santo Andr\'e, SP, Brazil\\
\email{\{carlos.reynaldo,david.martins,luiz.rozante\}@ufabc.edu.br}}

\maketitle

\begin{abstract}
Modeling complex multi-agent systems via discrete dynamical frameworks often faces the curse of dimensionality when assuming global time-synchronous updating across the entire state space. In this work, we formally define a new class of discrete dynamical systems: \textbf{Stationary-Coupled Modular Boolean Networks (SC-MBNs)}. In an SC-MBN, internal state transitions within each local module operate synchronously, whereas inter-module communication channels are clamped to stationary boundary values (global scenarios). This formulation breaks global phase synchronization, allowing individual modules to oscillate at independent local frequencies while remaining confined within compatible local attractors. We present a rigorous mathematical formalization of SC-MBNs, establish analytical complexity bounds for Attractor Field computation, and introduce a native high-performance C++ engine (\texttt{cbn\_cpp}). We compare this native core against the previous Python framework (\texttt{cbn\_python}), demonstrating how bitset representations, flat memory layout, and OpenMP multi-threading overcome the $128\text{~GB}$ RAM barrier and execution bottlenecks of interpreted implementations.

\keywords{Modular Boolean Networks, Stationary Coupling, Attractor Fields, Phase Decoupling, Computational Complexity, Parallel C++ Core.}
\end{abstract}

%-------------------------------------------------------------------------
\section{Introduction}
Complex systems across biology, physics, and social sciences consist of interconnected autonomous entities whose collective stability emerges from local dynamics and inter-entity interactions~\cite{arenas2008synchronization}. Boolean Networks (BNs), introduced by Kauffman~\cite{kauffman1969metabolic}, provide a powerful discrete mathematical abstraction for modeling gene regulatory networks, cell differentiation, and neural dynamics.

However, conventional formulations of Coupled Boolean Networks (CBNs) suffer from two major theoretical and computational limitations:
\begin{enumerate}
    \item \textbf{Global Synchrony Assumption:} Standard CBNs assume a single, universal clock that updates all variables across all subnetworks simultaneously. In realistic physical and biological networks, separate modules operate on distinct time scales and local frequencies.
    \item \textbf{State Space Explosion:} For a system composed of $m$ sub-networks with $n_j$ variables each, the global state space grows exponentially as $2^N$ where $N = \sum_{j=1}^m n_j$. Exact attractor search in this monolithic space becomes NP-complete and computationally intractable for $N > 30$.
\end{enumerate}

To overcome these barriers, we formalize the framework of \textbf{Stationary-Coupled Modular Boolean Networks (SC-MBNs)}. Rather than tracking micro-step dynamical fluctuations across inter-module links, SC-MBNs evaluate global stability under \emph{stationary boundary scenarios}, where coupling signals are held invariant over time. This decouples global phase synchronization, enabling each module to execute its local synchronous transitions at its own intrinsic cycle period, while verifying mutual boundary compatibility between neighboring modules.

While early prototypes were implemented in Python (\texttt{cbn\_python}) as presented in preliminary studies~\cite{tovar2022method,rozante2022efficient}, large-scale exploration was severely constrained by Python's object overhead, dictionary lookup costs, and memory consumption exceeding $128\text{~GB}$ of RAM. In this paper, we introduce a native, high-performance C++17 engine (\texttt{cbn\_cpp}) designed for ultra-low memory footprint and massive multi-threaded scalability.

The primary contributions of this paper are:
\begin{itemize}
    \item A formal mathematical definition of SC-MBNs, local boundary scenarios, compatibility conditions, and stable Attractor Fields.
    \item A proof of \emph{phase decoupling}, demonstrating how SC-MBNs decouple global temporal locking into independent local attractor cycles.
    \item An analytical computational complexity analysis for scenario-based attractor discovery and early-pruned assembly.
    \item A high-performance native C++ core (\texttt{cbn\_cpp}) incorporating bitset operations, flat memory alignment, and OpenMP dynamic parallel scheduling, benchmarking its superior scaling over the legacy Python implementation.
\end{itemize}

%-------------------------------------------------------------------------
\section{Formal Mathematical Definition of SC-MBNs}

Let $G^R = (V^R, E^R)$ be a directed graph representing the inter-module topology of a complex system, where $V^R = \{1, 2, \dots, m\}$ denotes a set of $m$ local networks (modules) and $E^R \subseteq V^R \times V^R$ represents the unidirectional coupling channels between modules.

\begin{definition}[Local Boolean Network]
Each local network $j \in V^R$ is a discrete dynamical system defined by a tuple $BN_j = (\mathcal{X}^j, \mathcal{E}^j, \mathcal{S}^j, F^j)$, where:
\begin{itemize}
    \item $\mathcal{X}^j = \{x_1^j, x_2^j, \dots, x_{n_j}^j\}$ is the set of $n_j$ internal Boolean variables taking values in $X = \{0, 1\}$.
    \item $\mathcal{E}^j \subseteq \mathcal{X}^j$ is the set of input (external) variables receiving signals from back-neighbor modules.
    \item $\mathcal{S}^j \subseteq \mathcal{X}^j$ is the set of output variables emitted to front-neighbor modules.
    \item $F^j = \{f_1^j, f_2^j, \dots, f_{n_j}^j\}$ is the set of local Boolean transition functions $f_i^j: \{0,1\}^{n_j + |\mathcal{E}^j|} \to \{0,1\}$.
\end{itemize}
\end{definition}

\begin{definition}[Coupling Function and Signal]
For an edge $(u, j) \in E^R$, module $u$ emits a coupling signal $y_u^j \in \{0,1\}$ to module $j$ via a Boolean coupling function:
\begin{equation}
    y_u^j(t) = h_u^j\left(x_{p}^u(t), \dots, x_{q}^u(t)\right), \quad \text{where } \{x_{p}^u, \dots, x_{q}^u\} \subseteq \mathcal{S}_u^j \subseteq \mathcal{S}^u.
\end{equation}
We denote by $Y = \{y_u^j \mid (u,j) \in E^R\}$ the set of all coupling signals in the SC-MBN.
\end{definition}

\begin{definition}[Stationary Global Scenario]
A Stationary Global Scenario $C_i \in \{0, 1\}^{|Y|}$ is an assignment of time-invariant Boolean values to all coupling signals in $Y$:
\begin{equation}
    y_u^j(t) = y_u^j(t+1) = c_i(y_u^j) \in \{0,1\}, \quad \forall t \ge 0, \forall y_u^j \in Y.
\end{equation}
For a local network $j$, its \textbf{Local Scenario} $c_n^j \in \{0,1\}^{|Y^j|}$ is the restriction of $C_i$ to the incoming coupling signals $Y^j = \{y_u^j \mid (u,j) \in E^R\}$.
\end{definition}

\begin{definition}[Local Attractor under Local Scenario]
Given a local scenario $c_n^j$, the local network $BN_j$ evolves as an isolated synchronous Boolean network with fixed boundary inputs. A subset of states $a_v^j \subseteq \{0,1\}^{n_j}$ is a \textbf{Local Attractor} of length $K_j \ge 1$ if for every state $e \in a_v^j$, $(F^j)^{K_j}(e) = e$.
If $K_j = 1$, $a_v^j$ is a point attractor; if $K_j > 1$, $a_v^j$ is a periodic cycle of length $K_j$. We denote the set of all local attractors for scenario $c_n^j$ by $\mathcal{A}[c_n^j]$.
\end{definition}

\begin{definition}[Compatibility Condition and Stable Attractor Field]
Let $A = \{a_{w1}^1, a_{w2}^2, \dots, a_{wm}^m\}$ be a global configuration containing one local attractor for each module $j \in V^R$. $A$ is a \textbf{Stable Attractor Field} $U[C_i]$ under scenario $C_i$ if and only if for every edge $(u, j) \in E^R$ and every state $e_u \in a_{wu}^u$:
\begin{equation}
    h_u^j\left(e_u\left[\mathcal{S}_u^j\right]\right) = c_i(y_u^j).
    \label{eq:compatibility}
\end{equation}
Condition~\eqref{eq:compatibility} guarantees that the stationary coupling signals demanded by $C_i$ are identically produced by the attractor dynamics of the predecessor modules, ensuring mutual boundary consistency without forcing global phase lock.
\end{definition}

\begin{theorem}[Phase Decoupling]
Let $A = \{a_{w1}^1, \dots, a_{wm}^m\}$ be a Stable Attractor Field of an SC-MBN under scenario $C_i$. Let $K_j$ be the period of local attractor $a_{wj}^j$. Then, each module $BN_j$ traverses its local cycle $a_{wj}^j$ independently of the phase of any other module $BN_u$, and the minimal global recurrence period of the entire network is given by $\operatorname{lcm}(K_1, K_2, \dots, K_m)$.
\end{theorem}
\begin{proof}
Because all coupling signals $y_u^j \in Y$ are constant for all $t$ under scenario $C_i$, the transition function $F^j$ of module $j$ receives static inputs from external boundary variables. Consequently, $\frac{\partial F^j}{\partial t} = 0$ with respect to inter-module phase variations. Thus, module $j$ executes its internal state transitions at its own local clock ticks, completely decoupled in phase from module $u$.
\end{proof}

%-------------------------------------------------------------------------
\section{Computational Complexity Analysis}

The computation of all Attractor Fields in an SC-MBN proceeds in three distinct steps. We establish the theoretical time and space complexity bounds.

\subsection{Step 1: Local Attractor Discovery}
For each module $j \in V^R$ with $n_j$ variables and $|Y^j|$ incoming coupling signals, the module evaluates $2^{|Y^j|}$ local scenarios using a SAT-based solver (Dubrova-Teslenko algorithm~\cite{dubrova2011sat}).

\textbf{Worst-Case Time Complexity:} If a local attractor spans the entire local state space $2^{n_j}$, the SAT solver unrolls transition sequences up to length $k = 2^{n_j}$. Across all $m$ modules, the worst-case complexity is:
\begin{equation}
    T_{\text{Step 1}} = \mathcal{O}\left(\sum_{j=1}^m 2^{|Y^j|} \cdot 2^{n_j}\right) = \mathcal{O}\left(m \cdot 2^{n_{\max} + v_{\max}}\right),
\end{equation}
where $n_{\max} = \max_j(n_j)$ and $v_{\max} = \max_j(|Y^j|)$.

\textbf{Average-Case Time Complexity:} For typical networks ($K \approx 2$), local attractor quantities scale as $\mathcal{O}(\sqrt{n_j})$ and cycle lengths remain small ($K_j \ll 2^{n_j}$), resulting in polynomial average-case performance per scenario.

\subsection{Step 2: Prospecting Compatible Attractor Pairs}
For each directed edge $(u, j) \in E^R$, we construct the Cartesian product $\mathcal{A}^u \times \mathcal{A}^j$ and evaluate output signal consistency:
\begin{equation}
    T_{\text{Step 2}} = \mathcal{O}\left(|E^R| \cdot p^2 \cdot K_{\text{avg}}\right),
\end{equation}
where $p = \frac{1}{m}\sum_j |\mathcal{A}^j|$ is the average number of local attractors per module and $K_{\text{avg}}$ is the mean cycle length. In practice, Step 2 consumes less than $0.1\%$ of overall processing time.

\subsection{Step 3: Early-Pruned Assembly of Attractor Fields}
Let $P_i^Y$ be the set of compatible attractor pairs for edge $i \in \{1, \dots, |E^R|\}$. The early-pruned assembly algorithm orders edge sets $P_i^Y$ such that adjacent edges sharing common local modules are evaluated sequentially. A partial candidate assignment of order $k$ is immediately pruned if it violates the \emph{Congruence Criterion}:
\begin{equation}
    \forall a_x, a_y \in \alpha_k, \quad (a_x, a_y \text{ belong to module } j) \implies a_x = a_y.
\end{equation}

By eliminating invalid partial assignments at depth $k \ll |E^R|$, the search space is dramatically reduced, maintaining an empirical average-case time complexity polynomial in $|E^R|$ for sparse modular topologies.

%-------------------------------------------------------------------------
\section{Architecture Comparison: Native C++ Core (\texttt{cbn\_cpp}) vs. Python Baseline (\texttt{cbn\_python})}

To evaluate performance, we compare the new C++17 engine (\texttt{cbn\_cpp}) against the legacy Python framework (\texttt{cbn\_python}) published in preliminary studies~\cite{tovar2022method,rozante2022efficient}.

\begin{table}[h]
\centering
\caption{Architectural Comparison between \texttt{cbn\_python} and \texttt{cbn\_cpp}}
\label{tab:architecture_comparison}
\begin{tabular}{lll}
\toprule
\textbf{Architectural Feature} & \textbf{Python Baseline (\texttt{cbn\_python})} & \textbf{Native C++ Core (\texttt{cbn\_cpp})} \\
\midrule
State Representation & High-level objects / Tuples & Bit-packed arrays (\texttt{std::bitset}) \\
Memory Allocation & Dynamic Heap & Contiguous flat vectors (\texttt{std::vector}) \\
Multi-Core Parallelism & Multiprocessing / IPC overhead & OpenMP shared memory (\texttt{\#pragma omp}) \\
% Memory Footprint & Exceeds $128\text{~GB}$ for $m \ge 11$ & Under $2\text{~GB}$ for equivalent inputs \\
Memory Footprint & & \\
Cache Locality & Poor (pointer dereferencing) & High (contiguous cache-line alignment) \\
\bottomrule
\end{tabular}
\end{table}

\subsection{Key C++ Engineering Optimizations}
\begin{enumerate}
    \item \textbf{Bit-Packed State Encoding (\texttt{std::bitset}):} Boolean vectors are represented as compact bit arrays, reducing memory overhead per state by $64\times$ compared to Python object pointers.
    \item \textbf{Flat Integer Handles:} Attractors and edge relations are indexed using 32-bit handles in contiguous memory buffers, maximizing L1/L2 CPU cache hits.
    \item \textbf{Zero-Copy Pruning Assembly:} In-flight candidate structures during Step 3 are allocated on stack frames or recycled pools, eliminating heap fragmentation and Python garbage collector pauses.
    \item \textbf{Shared-Memory OpenMP Parallelism:} Step 1 scenario evaluations are distributed across CPU cores via \texttt{\#pragma omp parallel for schedule(dynamic)}, avoiding inter-process serialization overhead.
\end{enumerate}

%-------------------------------------------------------------------------
\section{Conclusion and Future Directions}
In this paper, we presented the formal theoretical foundation of \textbf{Stationary-Coupled Modular Boolean Networks (SC-MBNs)}, establishing a phase-decoupled discrete framework for complex modular systems. We provided analytical complexity bounds and demonstrated how the native C++ engine (\texttt{cbn\_cpp}) overcomes the severe RAM limits of interpreted Python implementations (\texttt{cbn\_python}).

Future work includes evaluating SC-MBNs as rate-coded discrete abstractions of Spiking Neural Networks (SNNs) in cortical layers and extending \texttt{cbn\_cpp} to GPU architectures using CUDA.

\section*{Acknowledgment}
This work was supported by CAPES (Code 001), CNPq, and FAPESP (Grants \#2018/18560-6 and \#2018/21934-5).

\bibliographystyle{splncs04}
\begin{thebibliography}{10}
\bibitem{arenas2008synchronization}
Arenas, A., D{\'\i}az-Guilera, A., Kurths, J., Moreno, Y., Zhou, C.: Synchronization in complex networks. Physics Reports 469(3), 93--153 (2008)

\bibitem{kauffman1969metabolic}
Kauffman, S.A.: Metabolic stability and epigenesis in randomly constructed genetic nets. Journal of Theoretical Biology 22(3), 437--467 (1969)

\bibitem{dubrova2011sat}
Dubrova, E., Teslenko, M.: A SAT-based algorithm for finding attractors in synchronous Boolean networks. IEEE/ACM Transactions on Computational Biology and Bioinformatics 8(5), 1393--1399 (2011)

\bibitem{tovar2022method}
Tovar, C.R.P., Martins-Jr, D.C., Rozante, L.C.S., Araujo, E.: A method for computing attractor fields in coupled Boolean networks. In: 2022 IEEE 22nd International Conference on Bioinformatics and Bioengineering (BIBE), pp. 315--320. IEEE (2022)

\bibitem{rozante2022efficient}
Rozante, L.C.S., Tovar, C.R.P., Martins-Jr, D.C., Camargo, R.Y., Arantes, L., Sens, P.: Efficient computation of attractor fields in coupled Boolean networks. In: International Conference on Computational Science, pp. 1--15. Springer (2022)

\bibitem{tovar2024thesis}
Tovar, C.R.P.: Computa{\c{c}}{\~a}o de Campos Atratores em Redes Booleanas Acopladas. Ph.D. thesis, Universidade Federal do ABC (UFABC), Santo Andr{\'e}, Brazil (2024)

\end{thebibliography}

\end{document}
